% Source: https://artofproblemsolving.com/community/c3h1608486_happy_half_tau_day
% Stable identity: aops-topic-1608486-q01
A circle of diameter $d$ has area $\pi\left(\dfrac {d}{2}\right)^2=\dfrac{\pi d^2}{4}$. So the areas add exactly when the squares of the diameters add: the new diameter $D$ satisfies
\[
\dfrac{\pi D^2}{4}=\dfrac{\pi}{4}\left(12^2+13^2+14^2+15^2\right)\qquad\Longrightarrow\qquad D^2=144+169+196+225=734 .
\]

\begin{asy}
usepackage("amsmath");
size(12cm);
defaultpen(fontsize(10pt));
real[] d={12,13,14,15,sqrt(734)};
real x=0;
for(int i=0;i<5;++i){
  real r=d[i]/2; pair O=(x+r,0);
  draw(circle(O,r),linewidth(0.9));
  draw(O-(r,0)--O+(r,0),dashed);
  string s=(i<4)?"$"+string(d[i])+"$":"$\sqrt{734}$";
  label(s,O+(0,0.3*r));
  x+=d[i]+(i==3?6:2.5);
}
label("$+$",(12+1.25,0)+(0,-9));label("$+$",(12+2.5+13+1.25,0)+(0,-9));label("$+$",(12+2.5+13+2.5+14+1.25,0)+(0,-9));
label("$=$",(12+13+14+15+7.5+3,0)+(0,-9));
label("\mbox{}",(min(currentpicture,user=true)+max(currentpicture,user=true))/2,black); // keeps the text after the figure black when LaTeX typesets the labels inline
\end{asy}

Hence the new pie has diameter
\[
D=\boxed{\sqrt{734}\text{ inches}}\approx27.1\text{ inches}.
\]
(The dashed segments are the diameters; the figure is to scale.)
